\scanpage{95}
\mspage{94}
\margin{-7-}

consequences of 1.9, of which only the following (cf. 1.7) will be stated here, the paper comes to a close.

1.10 Theorem. \uline{Let} $k$ \uline{be a perfect field of} $\dim\leq1$ \uline{and} $G$ \uline{a connected linear group over} $k$. \uline{Then every conjugacy class over} $k$ \uline{contains an element over} $k$.

After the remark that Kneser, using extensions of 1.8, has recently shown (cf. 1.9) that $H^1(k,G)=0$ if $k$ is a p-adic field and $G$ a simply connected semisimple group over $k$, this introduction comes to a close.

\S2. \uline{Some recollections}. In this section we recall some known facts, including some characterizations 2.11 of regular semisimple elements, and establish some notations which are frequently used in the paper. If $k$ is a field, $k^*$ is its multiplicative group. The term ``algebraic group'' is often abbreviated to ``group''. If $G$ is a group, $G_{\unsure{0|O}}$ denotes its identity component. If $x$ is an element of $G$, then $G_x$ denotes the centralizer of $x$ in $G$, and $x_s$ and $x_u$ denote the semisimple and unipotent parts of $x$ when $G$ is linear. Assume now that $G$ is a semisimple group, that is, $G$ is a connected linear group with no nontrivial connected solvable normal subgroup. We write $r$ for the rank of $G$. Assume further that $T$ is a maximal torus in $G$ and that an ordering of the (discrete) character group of $T$ has been chosen. We write $\Sigma$ for the system of roots relative to $T$ and $X_\alpha$ for the subgroup corresponding to the root $\alpha$.

2.1 $X_\alpha$ \uline{is unipotent and isomorphic (as an algebraic group) to the}
